\begin{enumerate}
	\item On a $f(x) = (x+1)\,\text{e}^{x} = x\,\text{e}^{x} + \,\text{e}^{x}$. Donc les primitives de $f$ sont de la forme $F(x) = x\,\text{e}^{x} + c$ avec $c \in \mathbb{R}$. Donc $F(x) = 1 + x\,\text{e}^{x}$ est une primitive de $f$.
	
	\textbf{Réponse a.}
	\item $\vect{u_1} \begin{pmatrix}1\\1\\-1\end{pmatrix}$ est un vecteur directeur de $\left(\text{d}_1\right)$ et $\vect{u_2} \begin{pmatrix}-1\\1\\-1\end{pmatrix}$ est un vecteur directeur de $\left(\text{d}_2\right)$.
	
On a $\dfrac{1}{-1} \neq \dfrac{1}{1}$, donc $\vect{u_1}$ et $\vect{u_2}$ ne sont pas colinéaires. Les droites ne sont donc pas parallèles.
	
Supposons qu'elle soient sécantes en un point A. Alors les coordonnées $(a~;~b~;~c)$ de A vérifient les deux représentations paramétriques :
	
$\left\{\begin{array}{rcl}a &=& 2 + r \\ b &=& 1 + r \\ c &=& -r\end{array}\right.$ et $\left\{\begin{array}{rcl}a &=& 1 - s \\ b &=& -1 + s \\ c &=& 2 - s\end{array}\right.$.
	
On a donc $-r = 2 - s \iff r = s - 2$ et $2 + r = 1 - s \iff 2 + s - 2 = 1 - s \iff 2s = 1 \iff s = \dfrac{1}{2}$.
	
Donc : $b = -1 + \dfrac{1}{2} = -\dfrac{1}{2}$, $a = 1 - \dfrac{1}{2} = \dfrac{1}{2}$ et $c = 2 - \dfrac{1}{2} = \dfrac{3}{2}$.
	
Les deux droites sont sécantes en $\text{A}\left(\dfrac{1}{2}~;~-\dfrac{1}{2}~;~\dfrac{3}{2}\right)$.
	
	\textbf{Réponse a.}
	\item Un vecteur normal au plan (P) est $\vect{n} \begin{pmatrix}2\\-1\\2\end{pmatrix}$ et un vecteur directeur de $(\Delta)$ est $\vect{u} = \begin{pmatrix}1\\1\\-1\end{pmatrix}$.
	
	On a : $\vect{n}\cdot\vec{u} = 2 \times 1 - 1 \times 1 + 1 \times (-1) = 0$, donc $\vect{n}$ et $\vect{u}$ sont orthogonaux. Donc (P) et $(\Delta)$ sont parallèles.
	
	S'il existe un point commun à (P) et $(\Delta)$, alors $(\Delta)$ est incluse dans (P).
	
	$(\Delta)$ passe par le point de coordonnées $(2~;~4~;~1)$. Or $2 \times 2 - 4 + 1 - 1 = 0$, donc ce point est un point du plan (P). La droite $(\Delta)$ est donc incluse dans le plan (P).
	
	\textbf{Réponse b.}
	\item Un vecteur normal à $\left(\text{P}_1\right)$ est $\vect{n_1} = \begin{pmatrix}1\\-2\\1\end{pmatrix}$ et un vecteur normal à $\left(\text{P}_2\right)$ est $\vect{n_2} = \begin{pmatrix}2\\1\\1\end{pmatrix}$.
	
	On a $\dfrac{1}{2} \neq \dfrac{-2}{1}$, donc $\vect{n_1}$ et $\vect{n_2}$ ne sont pas colinéaires. Les plans ne sont donc pas parallèles. Ils sont donc sécants.
	
	De plus $\vect{n_1} \cdot \vect{n_2} = 1 \times 2 - 2 \times 1 + 1 \times 1 = 1$, donc les vecteurs $\vect{n}_1$ et $\vect{n}_2$ ne sont pas orthogonaux. Les plans ne sont donc pas perpendiculaires.
	
	\textbf{Réponse c.}
	\item On sait que $\vect{EF} \cdot \vect{EG} = \left\|\vect{EF}\right\| \times \left\|\vect{EG}\right\| \times \cos \left(\vect{EF},~\vect{EG}\right) \iff \cos \alpha = \dfrac{\vect{EF} \cdot \vect{EG}}{\left\|\vect{EF}\right\| \times \left\|\vect{EG}\right\|}$
	
On a : $\vect{EF}\begin{pmatrix}2-1\\4-2\\3-1\end{pmatrix} = \begin{pmatrix}1\\2\\2\end{pmatrix}$ et $\vect{EG}\begin{pmatrix}-2-1\\2-2\\5-1\end{pmatrix} = \begin{pmatrix}-3\\0\\4\end{pmatrix}$.
	$\vect{EF} \cdot \vect{EG} = 1 \times (-3) + 2 \times 0 + 2 \times 4 = 5$
	
	$\left\|\vect{EF}\right\| = \sqrt{1^2 + 2^2 + 2^2} = \sqrt{9} = 3$
	
	$\left\|\vect{EG}\right\| = \sqrt{(-3)^2 + 0^2 + 4^2} = \sqrt{25} = 5$
	
	Donc :
	$\cos \alpha = \dfrac{5}{3 \times 5} = \dfrac{1}{3}$, soit $\alpha \approx 71\degre$ d'après la calculatrice.
	\textbf{Réponse d.}
\end{enumerate}

\bigskip
